\subsection{COMBINATORIAL ANALYSIS}

PROPOSITION 9. Let E and F be two sets, a and b their cardinals, f a surjection of E onto F such that the sets \(\overline{f}^{1}(y)\) , where \(y \in F\) , all have the same cardinal c.~Then a = bc.

For the family \((\bar{f}^1 (y))_{\gamma \in \mathbb{F}}\) is a partition of \(\mathbf{E}\) , each element of which partition is a set of cardinal \(\mathfrak{c}\) ; hence the result (§ 3, no. 4, Proposition 6, Corollary 2).

DEFINITION 2. Let \(n\) be an integer. The product \(\prod_{i < n} (i + 1)\) is denoted by \(n!\) (read ``factorial \(n\)'').

We have \(0! = 1\) (Chapter II, §5, no. 3) and \(1! = 1\) . For each integer \(n\) we have \((n + 1)! = n!(n + 1)\) . This relation, together with the relation \(0! = 1\) , characterizes the term \(n!\) , as is easily seen by induction on \(n\) .

PROPOSITION 10. Let \(m\) and \(n\) be integers such that \(m \leqslant n\) . Then \(n! / (n - m)!\) is the number of injective mappings of a set \(A\) with \(m\) elements into a set \(B\) with \(n\) elements.

The proof is by induction on the number \(m \leqslant n\) of elements of A. If \(m = 0\) , the result is evident. Suppose that \(m + 1 \leqslant n\) . Let A be a set with \(m + 1\) elements, let \(A'\) be a subset of A with \(m\) elements, and let \(\{a\} = A - A'\) . Let F, \(F'\) be the sets of injective mappings of A, \(A'\) respectively into B, and let \(\varphi\) be the mapping \(f \to f|A'\) which maps each function \(f \in F\) to its restriction to \(A'\) . For each \(f' \in F'\) an element \(f\) of \(\overline{\varphi}(f')\) is uniquely determined by its value \(f(a)\) ; since \(f\) is injective, we must have \(f(a) \in B - f'(A')\) . It follows that \(\overline{\varphi}(f')\) has the same number \(n - m\) of elements as \(B - f'(A')\) . Hence, by Proposition 9, F has

\[
(n - m) \frac {n !}{(n - m) !} = \frac {n !}{(n - m - 1) !}
\]

elements by virtue of the inductive hypothesis. This completes the proof.

COROLLARY. The number of permutations of a finite set with n elements is equal to n!.

For this number is equal to the number of injections of the set into itself (§ 4, no. 2, Proposition 2, Corollary 4).

PROPOSITION 11. Let E be a finite set with n elements, and let \((p_i)_{1 \leqslant i \leqslant h}\) be a finite sequence of integers such that \(\sum_{i=1}^{h} p_i = n\) . Then the number of coverings \((\mathbf{X}_i)_{1 \leqslant i \leqslant h}\) of E by mutually disjoint sets \(\mathbf{X}_i\) such that \(\operatorname{Card}(\mathbf{X}_i) = p_i\) for \(1 \leqslant i \leqslant h\) is equal to

\[
n! / \prod_ {i = 1} ^ {h} p _ {i}!.
\]

Let G be the set of permutations of E and let P be the set of coverings \((\mathbf{X}_{i})_{1\leqslant i\leqslant h}\) which satisfy the conditions of the Proposition. Since \(\sum_{i=1}^{h}p_{i}=n\) , P is not empty. Let \((\mathbf{A}_{i})_{1\leqslant i\leqslant h}\) be an element of P. For each permutation \(f\in G\) the family \((f(\mathbf{A}_{i}))_{1\leqslant i\leqslant h}\) again belongs to P; let us denote it by \(\varphi(f)\) . For each element \((\mathbf{X}_{i})_{1\leqslant i\leqslant h}\) let us calculate the number of permutations \(f\in G\) such that \(\varphi(f)=(\mathbf{X}_{i})\) . We shall have \(\varphi(f)=(\mathbf{X}_{i})\) if and only if for each index i we have \(f(\mathbf{A}_{i})=\mathbf{X}_{i}\) . Hence the set of permutations f under consideration is equipotent to the product of the sets of bijections of \(A_{i}\) onto \(X_{i}\) (Chapter II, §4, no. 7, Proposition 8); consequently the set \(\overline{\varphi}((\mathbf{X}_{i})_{1\leqslant i\leqslant h})\) has \(\prod_{i=1}^{h}p_{i}!\) elements (Corollary to Proposition 10). Since G has n! elements, the result now follows from Proposition 9.

COROLLARY 1. Let A be a set with n elements and let p be an integer \(\leqslant n\) . Then the number of subsets of A which have p elements is \(\frac{n!}{p!(n-p)!}\) . Put \(h=2, p_{1}=p, p_{2}=n-p\) in Proposition 11.

The number of subsets containing \(p\) elements in a set of \(n\) elements (where \(p \leqslant n\) ) is denoted by \(\binom{n}{p}\) and is called the binomial coefficient with indices \(n\) and \(p\) . From the relation \(\binom{n}{p} = \frac{n!}{p!(n - p)!}\) it follows immediately that \(\binom{n}{p} = \binom{n}{n - p}\) .

This is also a consequence of the fact that, if E is a set with n elements, \(X \rightarrow E - X\) is a bijection of the set of subsets of E consisting of p elements onto the set of subsets of n - p elements.

We put \(\binom{n}{p} = 0\) for each pair of natural integers such that \(p > n\) . With this convention the number of subsets of \(p\) elements in a set of \(n\) elements is \(\binom{n}{p}\) for every natural integer \(p\) .

COROLLARY 2. Let E and F be totally ordered finite sets with p and n elements, respectively. Then the number of strictly increasing mappings of E into F is \(\binom{n}{p}\) .

For such a mapping is an injection of E into F (§1, no. 12, Proposition 11), and since E and F are well-ordered (§4, no. 4, Corollary 1 to Proposition 3), for each subset X of p elements of F there is exactly one strictly increasing mapping of E onto X (§2, no. 5, Theorem 3).

PROPOSITION 12. For each integer \(n\) , we have

\[
\sum_ {p} \binom{n}{p} = 2 ^ {n}.
\]

For if E is a set of n elements, the left-hand side of the equality is the number of subsets of E. Now apply Proposition 12 of § 3, no. 5.

PROPOSITION 13. If n and p are integers, then

\[
\binom {n + 1} {p + 1} = \binom {n} {p + 1} + \binom {n} {p}.
\]

Let E be a set with \(n + 1\) elements, let P be the set of all subsets of E containing \(p + 1\) elements, let a be an element of E, and put

\[
\mathbf {E} ^ {\prime} = \mathbf {E} - \{a \}.
\]

Let \(P'\) (resp. \(P''\) ) denote the set of subsets of \(p + 1\) elements of E which contain a (resp. do not contain a). The set \(P''\) is the set of subsets of \(p + 1\) elements of \(E'\) and therefore has \(\binom{n}{p+1}\) elements. The mapping \(X \to X \cap E'\) is a bijection of \(P'\) onto the set of subsets of p elements of \(E'\) , and \(P'\) therefore has \(\binom{n}{p}\) elements. The result now follows from the fact that P is the union of the disjoint sets \(P'\) and \(P''\) .

Proposition 13 can also be proved by means of a simple calculation from the formula \(\binom{n}{p} = \frac{n!}{p!(n - p)!}\) for \(p \leqslant n\) .

PROPOSITION 14. Let \(n\) be an integer \(> 0\) . Then the number \(a_{n}\) (resp. \(b_{n}\) ) of ordered pairs \((i, j)\) of integers such that \(1 \leqslant i \leqslant j \leqslant n\) (resp. \(1 \leqslant i < j \leqslant n\) ) is \(\frac{1}{2} n(n + 1)\) ((resp. \(\frac{1}{2} n(n - 1)\) ).

For \(b_{n}\) is the number of subsets of 2 elements in \([1, n]\) ; hence

\[
b _ {n} = \frac {n !}{2 ! (n - 2) !} = \frac {1}{2} n (n - 1).
\]

The value of \(a_{n}\) is deduced from this by noting that the set of ordered pairs \((i, j)\) such that \(1 \leqslant i \leqslant j \leqslant n\) is the union of the set of ordered pairs \((i, j)\) such that \(1 \leqslant i \leqslant j < n\) and the set of pairs \((i, i)\) where \(1 \leqslant i \leqslant n\) . Thus \(a_{n} = n + b_{n} = \frac{1}{2} n(n + 1)\) .

COROLLARY. For each integer \(n > 0\) , we have

\[
\sum_ {i = 1} ^ {n} i = \frac {1}{2} n (n + 1).
\]

In the set A of ordered pairs of integers \((i, j)\) such that \(1 \leqslant i \leqslant j \leqslant n\) , let \(A_{k}\) denote the subset of pairs \((i, k)\) , where \(1 \leqslant i \leqslant k\) (for an arbitrary integer \(k \leqslant n\) ). Then \(A_{k}\) has k elements. But \((\mathrm{A}_{k})_{1 \leqslant k \leqslant n}\) is a partition of A; hence the result.

PROPOSITION 15. Let \(n\) and \(h\) be integers and let \(E\) be a set with \(h\) elements. Then the number of mappings \(u\) of \(E\) into \([0, n]\) such that \(\sum_{x \in E} u(x) \leqslant n\) (resp. \(\sum_{x \in E} u(x) = n\) , for \(h > 0\) ) is

\[
\binom {n + h} {h} \left(\text { resp. } \binom {n + h - 1} {h - 1}\right).
\]

Let \(\mathrm{A}(h, n)\) (resp. \(\mathrm{B}(h, n)\) ) denote the number of mappings \(u\) of \(\mathbf{E}\) into \([0, n]\) such that \(\sum_{x \in \mathbf{E}} u(x) \leqslant n\) (resp. \(\sum_{x \in \mathbf{E}} u(x) = n\) for \(h > 0\) ). We show first that \(\mathrm{A}(h - 1, n) = \mathrm{B}(h, n)\) . For this, let \(\mathbf{E}'\) be a subset of \(\mathbf{E}\) with \(h - 1\) elements, and let \(\{a\} = \mathbf{E} - \mathbf{E}'\) . If \(u\) is a mapping of \(\mathbf{E}\) into \([0, n]\) such that \(\sum_{x \in \mathbf{E}} u(x) = n\) , then its restriction \(u'\) to \(\mathbf{E}'\) is such that \(\sum_{x \in \mathbf{E}'} u'(x) \leqslant n\) , and moreover \(u(a) = n - \sum_{x \in \mathbf{E}'} u'(x)\) . Conversely, every mapping \(u'\) of \(\mathbf{E}'\) into \([0, n]\) which satisfies \(\sum_{x \in \mathbf{E}'} u'(x) \leqslant n\) defines a unique mapping \(u\) of \(\mathbf{E}\) into \([0, n]\) , of which \(u'\) is the restriction, and which is such that \(\sum_{x \in \mathbf{E}'} u(x) = n\) .

We note next that if \(\sum_{x\in E}u(x)\leqslant n\) , then either \(\sum_{x\in E}u(x) = n\) , or else \(\sum_{x\in E}u(x)\leqslant n - 1\) , and these two possibilities are mutually exclusive. Consequently

\[
\mathrm{A} (h, n) = \mathrm{A} (h, n - 1) + \mathrm{B} (h, n) = \mathrm{A} (h, n - 1) + \mathrm{A} (h - 1, n).
\]

Since \(\mathbf{A}(0,0) = 1 = \binom{0}{0}\) , the formula \(\mathbf{A}(h,n) = \binom{n + h}{h}\) follows from above and from Proposition 13, by induction on \(n + h\) .

* The number of monomials \(\mathbf{X}_1^{\alpha_1}\mathbf{X}_2^{\alpha_2}\ldots \mathbf{X}_h^{\alpha_n}\) in \(h\) indeterminates with total degree \(\leqslant n\) is evidently equal to the number of mappings \(i\to \alpha_{i}\) of \([1,h]\) into \([0,n]\) such that \(\sum_{i = 0}^{h}\alpha_{i}\leqslant n,\) and is therefore equal to \((n + h)\) by Proposition 15; this number is also the number of monomials in \(h + 1\) indeterminates with total degree \(n.\) *

